\documentclass[11pt]{article}
\usepackage[margin=0.82in]{geometry}
\usepackage{amsmath,amssymb,amsfonts}
\usepackage{booktabs,array,xcolor}
\usepackage{enumitem}
\usepackage{hyperref}
\hypersetup{colorlinks=true,linkcolor=blue,urlcolor=blue}
\newcommand{\R}{\mathbb{R}}
\newcommand{\C}{\mathbb{C}}
\newcommand{\E}{\mathbb{E}}
\newcommand{\softplus}{\operatorname{softplus}}
\newcommand{\softmax}{\operatorname{softmax}}
\newcommand{\LN}{\operatorname{LayerNorm}}
\newcommand{\Lam}{\Lambda}
\title{\textbf{S2P2 vs.\ SS2P2}\\[2pt]
\large Same expressive backbone, two output heads: how the gated bound turns a
state-space point process into a \emph{stable} simulator}
\author{}\date{}
\begin{document}\maketitle\vspace{-3em}

\section{One-line summary}
\noindent\textbf{S2P2} is a state-space point process: a stacked latent-linear-Hawkes (LLH)
backbone with a \emph{Proj\,\&\,Softplus} rate head and a rate-coupled mark head.
\textbf{SS2P2} keeps that backbone \emph{verbatim} and swaps only the two heads for a
\textbf{gated-bounded (G1) rate} and a \textbf{rate-neutral mark head}:
\[
\textbf{SS2P2}=\underbrace{\text{S2P2 LLH backbone}}_{\text{unchanged, expressive}}
\;\to\;\underbrace{\text{G1 bounded rate }\lambda(t)}_{\text{Poisson sandwich}}
\;\times\;\underbrace{\text{rate-neutral }p^*(x\mid t)}_{\text{decoupled marks}} .
\]
The backbone is identical, so any difference in behaviour is attributable to the heads.

\section{Shared backbone (both models)}
Stacked LLH/SSM layers $\ell=1,\dots,L$ with diagonal state, ZOH-discretized and
parallel-scan computable, with a paper-faithful LayerNorm'd readout producing the
embedding $u(t):=u^{(L)}(t)\in\R^{H}$:
\[
x^{(\ell)}_{t'}=\bar A^{(\ell)} x^{(\ell)}_{t}+(\bar A^{(\ell)}-I)\,\tilde B^{(\ell)} u^{(\ell-1)}
\;[+\,\tilde E^{(\ell)}\alpha\ \text{at an event}],\qquad
u^{(\ell)}=\LN\!\bigl(\mathrm{GELU}(y^{(\ell)})+u^{(\ell-1)}\bigr).
\]
Both models read the \emph{same} $u(t)$. They differ only in what comes next.

\section{The two heads, side by side}
\begin{center}\small\renewcommand{\arraystretch}{1.35}
\begin{tabular}{>{\raggedright\arraybackslash}p{2.4cm} >{\raggedright\arraybackslash}p{6.0cm} >{\raggedright\arraybackslash}p{6.0cm}}
\toprule
 & \textbf{S2P2} (baseline) & \textbf{SS2P2} (stable variant)\\
\midrule
Rate $\lambda(t)$ &
$\lambda=s\cdot\softplus(w^\top u+b)$ \newline (\emph{Proj\,\&\,Softplus}; unbounded above) &
$h=\sigma(W_o u)\odot\tanh(u)\in(-1,1)^H$, \newline
$\lambda=s\cdot\softplus(w_{\mathrm{eff}}^\top h+b)$, $\;\|w_{\mathrm{eff}}\|_1\le c$ \newline
(\textbf{gated, two-sided bounded})\\
\addlinespace
Mark $p^*(x\mid t)$ &
$p(k)=\lambda_k/\sum_j\lambda_j$ — \emph{coupled} to the rate field (or a shared-head softmax) &
$p^*(k)=\softmax(\mathrm{MLP}(u))_k$ — \emph{decoupled}, rate-neutral\\
\addlinespace
$\lambda_x^*(t)$ & emergent from the field & $\lambda(t)\,p^*(x\mid t)$, with $\sum_x p^*=1$\\
\addlinespace
Rate range & $(0,\infty)$ — can run away in closed loop & $\bigl(s\,\softplus(b-c),\,s\,\softplus(b+c)\bigr)$ — finite both sides\\
\addlinespace
Mean calibration & emergent (MLE) & emergent (MLE); $s,b$ set the level — \emph{no} EMA / branching pin\\
\bottomrule
\end{tabular}
\end{center}

\paragraph{Why the gate matters (G1).} Because $|h_i|<1$ and $\|w_{\mathrm{eff}}\|_1\le c$, we have
$|w_{\mathrm{eff}}^\top h|<c$, hence the closed-form two-sided bound
\[
\lambda(t)\in\bigl(\ell_-,\ell_+\bigr),\quad
\ell_\pm=s\,\softplus(b\pm c)
\;\;\Rightarrow\;\;
\mathrm{Poisson}(\ell_-)\preceq N(t)\preceq\mathrm{Poisson}(\ell_+).
\]
This is the ``Poisson sandwich'': in a free-running rollout the event rate can neither
explode (no runaway) nor die (no silence), \emph{without} any subcriticality projection.
S2P2's $\softplus(w^\top u+b)$ has no upper bound, so a self-exciting closed loop can
diverge. The magnitude lives in the learnable scale $s$ and the dynamic range in
$\softplus(b\pm c)$, so the bound is \emph{multiplicative} (wide) rather than additive.

\paragraph{Why decoupled marks.} S2P2's $p(k)=\lambda_k/\sum\lambda_j$ ties the mark
distribution to the rate field. SS2P2 reads marks from the same $u$ through an independent
softmax MLP, so marks stay fully expressive while never perturbing the bounded total rate
($\sum_x\lambda^*_x=\lambda$).

\section{Stability properties}
\begin{center}\small\renewcommand{\arraystretch}{1.2}
\begin{tabular}{lcc}
\toprule
property & S2P2 & SS2P2\\
\midrule
No runaway in closed-loop rollout & \textcolor{red}{not guaranteed} & \textcolor{teal}{\textbf{guaranteed}} ($\lambda\le\ell_+$)\\
No death / silent absorbing state & \textcolor{red}{not guaranteed} & \textcolor{teal}{\textbf{guaranteed}} ($\lambda\ge\ell_->0$)\\
Self-limiting gain & no & yes ($\tanh$ saturates, $\softplus'<1$)\\
Off-distribution fail-safe & no & yes ($o\!\to\!0\Rightarrow\lambda\!\to\!s\,\softplus(b)$)\\
Marks perturb the rate? & yes (coupled) & no (rate-neutral)\\
Expressive (parallel-scan) backbone & yes & yes (identical)\\
\bottomrule
\end{tabular}
\end{center}

\section{Empirical head-to-head}
Identical settings (seq 64 / stride 32, categorical mark head, same \texttt{genuine\_eval} +
\texttt{stylized\_facts} harness). \emph{Prediction}: top-1 mark accuracy / perplexity over
genuine events. \emph{Simulation}: relative error of stylized facts vs.\ real (lower is
better); SIM is their mean.

\begin{center}\small\renewcommand{\arraystretch}{1.25}
\begin{tabular}{lcccccc}
\toprule
model & ACC $\uparrow$ & PPL $\downarrow$ & Fano\_re & clus\_re & kurt\_re & \textbf{SIM} $\downarrow$\\
\midrule
\textbf{SS2P2} & 0.4519 & 8.51 & 0.939 & \textbf{0.241} & 0.658 & \textbf{0.613}\\
S2P2 (parent)  & \textbf{0.4648} & \textbf{8.47} & \textbf{0.388} & 1.331 & \textbf{0.159} & 0.626\\
\bottomrule
\end{tabular}
\end{center}

\noindent Raw stylized facts (real $\to$ model):
\begin{itemize}[leftmargin=1.4em,itemsep=1pt]
\item Volatility clustering ($|r|$-ACF, lags 1--10): real $0.127\to$ \textbf{SS2P2 $0.096$}, S2P2 $0.296$ (S2P2 \emph{over}-clusters).
\item Excess kurtosis: real $8.79\to$ SS2P2 $3.01$, \textbf{S2P2 $7.39$}.
\item Fano @1s: real $54.9\to$ SS2P2 $5.77$, \textbf{S2P2 $26.0$}.
\end{itemize}

\paragraph{Read.} Prediction is essentially tied (SS2P2 $-1.3$\,pp ACC, equal PPL): the
decoupled softmax head keeps S2P2's predictive power. Aggregate SIM is tied
($0.613$ vs $0.626$), but the \emph{profile} differs: the G1 bound regularizes the rate, so
SS2P2 matches volatility clustering far better while undershooting Fano/kurtosis (the bound
plus a tight $\|w\|_1$ cap limit burst magnitude — recoverable with a looser cap / more
training). The point SS2P2 buys is \textbf{stability}: the bounded rate ran a closed-loop
rollout with no blow-up, by construction.

\vspace{0.4em}
\noindent\footnotesize\textit{Caveat: these numbers are a local validation on Coinbase-BTC
(1 day), 10 epochs on CPU, short rollout — smoke-quality, for isolating the head change. The
paper-grade comparison (Gemini data, full config grid, GPU) is staged for the cluster.}

\end{document}
